% =====================================================================
%   摘要 (Abstract)
% =====================================================================

\section*{摘\quad 要}

近年来，大语言模型（Large Language Model, LLM）作为智能体的潜力已被广泛认可。
因此，迫切需要在具有挑战性的交互式环境任务上，对将 LLM 作为智能体的能力
进行定量评估。本文提出 \textbf{AgentBench}——一个由 8 个不同环境组成的多维度
基准（benchmark），用以评估 LLM 作为智能体（LLM-as-Agent）的推理与决策能力。
我们在 29 个基于 API 的商用 LLM 与开源（Open-Sourced, OSS）LLM 上进行了大规模
测试，结果显示：尽管顶尖商用 LLM 在复杂环境中展现出强大的智能体行为能力，
但其与众多规模不超过 70B 的开源竞品之间仍存在显著的性能差距。
我们进一步识别了环境与 LLM 各自的典型失败原因，发现较弱的长期推理、决策与
指令遵循能力是当前阻碍 LLM 智能体走向实用化的主要瓶颈。提升指令遵循能力，
并在高质量多轮对齐数据上进行训练，可有效改进智能体表现。此外，与既有
认知不同，本文发现代码训练对不同智能体任务的影响具有两面性。

本文相关的数据集、环境与一体化评估包已发布于
\url{https://github.com/THUDM/AgentBench}。

\begin{figure}[t]
    \centering
    \includegraphics[width=0.85\linewidth]{overview_placeholder.png}
    \caption{
        本图为原论文图 1 的中文译图——\textbf{LLM 在 AgentBench 上的总体表现}。
        左侧将各环境下的典型 LLM 表现与最佳模型对齐显示；右侧给出 29 个
        LLM 在 8 个环境上的总体得分，虚线分别表示 API 类商用 LLM 与
        OSS LLM 的平均分。可以看出：商用 LLM 平均得分显著高于 OSS LLM，
        但即便是最强的 \code{gpt-4}，距离真正实用化的智能体仍有可观差距。
    }
    \label{fig:overview}
\end{figure}

% 注释：保持章节连续，自然从 \section{引言} 开始即可（见 intro.tex）。
% 本翻译稿中不再重置计数器，避免后续页被推到下一页。
